% Generated by backend/scripts/make_figures.py. Do not edit; rerun the script.

\begin{figure}[htbp]\centering
\includegraphics[width=\textwidth]{figures/generated/fig01_class_balance.pdf}
\caption{\textbf{Class balance and base rate.} Labelled tokens alive at launch + 5 minutes (983), by how they were sampled (a) and the share labelled RUG (b). The case-control strata deliberately over-sample rare outcomes, so their RUG rate (29.0\%) is a design choice, not a property of Base; the random cohort gives the true rate, which moved from 7.9\% before the 20 Apr 2026 split to 45.0\% after it. Any metric read without the random cohort's base rate beside it is uninterpretable.}\label{fig:fig01_class_balance}
\end{figure}

\begin{figure}[htbp]\centering
\includegraphics[width=\textwidth]{figures/generated/fig02_label_sensitivity.pdf}
\caption{\textbf{Label sensitivity.} How many of the 1,058 labelled tokens change class when one labelling threshold is moved and the rest are held at their current values (open circle: the value in use; each panel has its own scale). Every verdict is re-derived from the stored pool observations. One step either side of the value in use moves at most 42 tokens (4.0\%). The label is not robust everywhere: calling a token dead only if it never trades again (dead-day cutoff 0) moves 648 tokens, and observing for fewer than 60 days instead of 90 moves 295-320, because FAILED needs a long observation to be told apart from a slow token.}\label{fig:fig02_label_sensitivity}
\end{figure}

\begin{figure}[htbp]\centering
\includegraphics[width=\textwidth]{figures/generated/fig03_launch_dates.pdf}
\caption{\textbf{Launch dates per class.} Weekly launch counts per class for every labelled token alive at minute 5, with the temporal split on 20 April 2026 (vertical line). Models are trained only on launches before it and tested on launches after it, so a model cannot learn from the period it is judged on. The class mix is not stable across the split: after it, RUG dominates, which is why results are reported per period.}\label{fig:fig03_launch_dates}
\end{figure}

\begin{figure}[htbp]\centering
\includegraphics[width=\textwidth]{figures/generated/fig04_pr_curves.pdf}
\caption{\textbf{Precision-recall, three baselines.} Precision-recall on the random cohort before the split (179 tokens, 64 buyer losses, base rate 0.36) for the target 'a buy at minute 5 is worth half or less at hour 6'. Models were trained on the case-control strata, with no operator family on both sides. Gradient boosting reaches PR-AUC 0.83 (95\% family bootstrap 0.73-0.91); the fixed rule checkers sit near the base rate (0.37). This set is out of sample but not out of time (Figure 6).}\label{fig:fig04_pr_curves}
\end{figure}

\begin{figure}[htbp]\centering
\includegraphics[width=\textwidth]{figures/generated/fig05_reliability.pdf}
\caption{\textbf{Reliability before and after calibration.} Reliability curves (points: equal-count bins; dashed line: perfect calibration). (a) On the random cohort before the split, raw boosting (Brier 0.127) against the same scores after Platt scaling fitted out of fold by operator family (Brier 0.129; base-rate Brier 0.230). (b) The frozen live model on launches it had never seen: its scores average 0.096 against an observed rate of 0.046, and its Brier (0.043) is worse than always predicting the base rate (0.044). The model ranks well forward but its probabilities are not calibrated, which is why the product shows the forward band, not the raw score.}\label{fig:fig05_reliability}
\end{figure}

\begin{figure}[htbp]\centering
\includegraphics[width=\textwidth]{figures/generated/fig06_decay.pdf}
\caption{\textbf{Performance by test period.} (a) A model trained only on launches before 1 March 2026, scored on each later month of the random cohort. On 1 March-17 April combined it reaches PR-AUC 0.796 (57 tokens; the early-concentration checker 0.331). From May, buyer losses almost vanish from the data, so PR-AUC rests on a handful of positives or is undefined. (b) The frozen live model, week by week, on launches that had not happened when it was trained; weeks with fewer than 50 resolved launches are omitted, the weeks around Base's block-time switch (5-16 November) are excluded by rule, and the last week is still resolving oldest-first. Performance moves with the base rate and the launch mix (a scam wave on one launchpad template from 3 November), not steadily down.}\label{fig:fig06_decay}
\end{figure}

\begin{figure}[htbp]\centering
\includegraphics[width=\textwidth]{figures/generated/fig07_importance.pdf}
\caption{\textbf{Permutation feature importance.} The 20 features whose shuffling most reduces gradient boosting's PR-AUC on the random cohort before the split (30 repeats each; bars ±1 s.d.). The top feature, price run-up by minute 5, costs 0.335. Graph features (g\_) carry little weight, consistent with Figure 8. Importance is measured on held-out data, so it says what the model relies on there, not what causes a loss.}\label{fig:fig07_importance}
\end{figure}

\begin{figure}[htbp]\centering
\includegraphics[width=\textwidth]{figures/generated/fig08_three_way.pdf}
\caption{\textbf{Timing, graph, or both.} PR-AUC with 95\% family-bootstrap intervals for each feature group on the random cohort before the split. With gradient boosting, timing features alone reach 0.83 and both together 0.83: the intervals overlap, so the wallet-graph features add nothing measurable over entry timing. Graph features alone reach 0.54. This is the outcome the project reports as its graph result (DECISIONS.md, Phase 5 Part D and Phase 6).}\label{fig:fig08_three_way}
\end{figure}

\begin{figure}[htbp]\centering
\includegraphics[width=\textwidth]{figures/generated/fig09_comovement.pdf}
\caption{\textbf{Co-movement lift.} How often two wallets trade the same token at about the same time: pairs inside one funding cluster, random pairs, and a degree-preserving shuffle of the cluster graph. Cluster pairs co-move 3.6$\times$ as often as random pairs when only buys count (p = 0.049) and 4.7$\times$ for buys or sells (p = 1e-05); the shuffled control falls to 0.26$\times$, so the effect is in the clusters, not the degree structure. The counts are small (3 of 40 cluster pairs for buys only), and the clusters are token-acquisition provenance, not ETH funding (Section on gaps).}\label{fig:fig09_comovement}
\end{figure}

\begin{figure}[htbp]\centering
\includegraphics[width=\textwidth]{figures/generated/fig10_cluster_case.pdf}
\caption{\textbf{A funding cluster, drawn.} A real cluster from the store (funder 0x46876d5b…, a contract). (a) One transaction sends a token to 81 wallets. (b) 41 minutes later, a single swap transaction buys token 0x778711… and delivers it to 81 of them; 6 minutes after that, another does the same with token 0xa397e1…. This is bundling: the wallets look like many buyers on a token's holder list, but they are one actor's 2-transaction operation. It is not 81 independent decisions, and the search for independent wallets buying together within 90 seconds found none in the store.}\label{fig:fig10_cluster_case}
\end{figure}

\begin{figure}[htbp]\centering
\includegraphics[width=\textwidth]{figures/generated/fig11_pipeline.pdf}
\caption{\textbf{Pipeline.} The system as built. Research path (top): keyless RPC reads of Base go through ingestion into PostgreSQL, are labelled by rule, and become 23 features at launch + 5 minutes for the models. Live path (bottom): the scorer judges every new pool once at minute 5 with the frozen model, the resolver measures the outcome at hour 6 and later horizons, and the append-only forward record and paper agent are built from those rows. The API serves both, plus proxied market data, to the React front end, where every section is badged live or sample. Durations are converted to blocks through one clock that knows Base's 2-second and 0.2-second regimes.}\label{fig:fig11_pipeline}
\end{figure}
